\documentclass[11pt]{article}
\usepackage[margin=1in]{geometry}
\usepackage[T1]{fontenc}
\usepackage{lmodern}
\usepackage{amsmath,amssymb,bm,booktabs,tabularx,array}
\usepackage{graphicx,float,placeins,microtype,fancyhdr}
\usepackage[hidelinks]{hyperref}
\usepackage{xurl}
\hypersetup{pdftitle={HPS GPR v6.1: reconstructed signal templates},pdfauthor={Emrys Peets}}
\pagestyle{fancy}\fancyhf{}\setlength{\headheight}{14pt}
\fancyhead[L]{\small HPS Gaussian-Process Resonance Search}
\fancyhead[R]{\small v6.1 / MC signal templates}\fancyfoot[C]{\thepage}
\fancypagestyle{plain}{\fancyhf{}\fancyhead[L]{\small HPS Gaussian-Process Resonance Search}
\fancyhead[R]{\small v6.1 / MC signal templates}\fancyfoot[C]{\thepage}}
\setlength{\emergencystretch}{2em}\setlength{\tabcolsep}{4pt}
\renewcommand{\arraystretch}{1.08}
\newcommand{\eps}{\varepsilon}
\newcommand{\fig}[3]{\begin{figure}[H]\centering
\includegraphics[width=.98\linewidth,height=#1\textheight,keepaspectratio]{#2}\caption{#3}\end{figure}}
\title{HPS Gaussian-Process Resonance Search\\[.5em]
\large Selected Reconstructed-MC Templates\\Standalone Study v6.1}
\author{Emrys Peets\\Stanford University and SLAC National Accelerator Laboratory}
\date{23 September 2026}
\begin{document}
\maketitle
\begin{abstract}
The 2021 Gaussian signal model is compared with normalized reconstructed
target-constrained selected-candidate mass distributions from twelve A-prime MC
masses between 40 and 260~MeV. The baseline retains the original $\pm2.25\sigma_m$ fit and
training mask, observed 10\% spectrum, GP kernel states and coupling
normalization. Additional common mask half-widths of 2.4, 2.5 and 2.6
nominal sigma are compared separately. Direct-template observed 90\% $CL_s$ limits
and local $p_0$ values at the generated masses are compared with Gaussian
fits to the same data. A dense morphing scan is reported separately as
exploratory. The MC production is v13 and the observed sample is v16;
their target constraint and momentum-sum labels agree, but timing
selections differ and full production equivalence is not established.
Stored reconstructed masses receive no added smearing or recentering.
The inherited coupling coordinate is a conditional comparison, not a
validated A-prime exclusion: an unmatched or accidental continuum need
not scale with the signal coupling.
A further comparison centers the primary fit and training window on an MC-only estimate of the reconstructed core, retaining the nominal width and unchanged MC shape.
\end{abstract}

\section{Observed comparison and scope}
The opening comparison retains the original pole-centered prescription.
The MC core-centered update appears in Section~\ref{sec:corecenter}.
At the 10 shared generated masses, the median MC-to-Gaussian limit ratio is 0.9876; the range is 0.5332--3.7590. The median absolute relative limit change is 22.05\%. The largest absolute relative limit change occurs at 60~MeV (+275.90\%). The smallest direct-MC local probability among these tested points occurs at 80~MeV, with $p_0$ = $2.477\times10^{-2}$ and $Z_0=1.964$. This is a sparse-grid local comparison, not a scan-global probability.


The direct comparison uses generated masses from 60 through 240~MeV in
20~MeV steps, inside the original 2021 search. The 40 and 260~MeV samples
provide shape and interpolation information outside those direct test
points. Full 2015 and full 2016 Gaussian references are retained in the
package, but no matching reconstructed MC for those campaigns was
supplied. The 2021 response is therefore not substituted into their fits.

The result is conditional on the supplied selected-MC shapes and their
finite statistics. It is neither a newly qualified detector-response
model nor a coverage calibration. The supplied candidate distributions
have broad components and substantial displacement from the generated
pole at several masses; they are not established as pure reconstructed
resonance responses. An observed difference between the two
curves can reflect mass bias, core width, non-Gaussian tails, selection
differences and finite-MC fluctuations together.

\clearpage
\section{Reconstructed mass distributions}
The source histograms retain selected target-constrained masses within
0--400~MeV and record out-of-range counts separately. The pole mass sets the hypothesis; the reconstructed mean
is allowed to differ. The comparison does not shift the MC peak onto the
pole mass or force its width to the Gaussian parameterization.

\fig{.76}{../figures/mc_shapes.pdf}{Unit-normalized stored reconstructed
MC mass densities and the nominal Gaussian at each generated mass.
Densities span the full stored range on a logarithmic axis so broad or
displaced populations remain visible.
Dotted lines mark the original $\pm2.25\sigma_m$ boundaries, calculated
from the nominal Gaussian resolution. Displayed densities are normalized
over their stored histogram range; fitted templates retain their
separately declared support normalization. No extra smearing or
recentering is applied.}

\clearpage
\section{Selected MC statistics and shape properties}
The source histograms store the weighted bin contents and squared-weight
sums. The effective number of weighted entries is
\begin{equation}
 N_{\mathrm{eff}}=\frac{(\sum_jw_j)^2}{\sum_jw_j^2}.
\end{equation}
It describes MC sampling, not data exposure or signal efficiency.
Selected rows need not be unique physical events unless the production
candidate policy establishes that identity. The source files, latest
ROOT tree cycles, branches, weights and extraction counts are recorded
in the histogram metadata and manifest.

The stored 0--400~MeV range contains 99.9848--99.9978\% of the selected weight across samples; overflow counts are shown separately. All extracted weights are unit weights, so $N_{\rm eff}$ equals the selected count. The raw particle-vector momentum diagnostic falls below 2.8~GeV for 1.35--8.28\% of rows. A bounded read finds that stored \texttt{psum} and \texttt{psum\_scalar} satisfy the cut in every inspected row; the raw-vector discrepancy therefore reflects an incompatible convention, not an established selection failure.


\begin{table}[H]\centering\small
\begin{tabular}{rrrrrr}
\toprule
$m_0$ (MeV) & Selected rows & Overflow & Median bias & $\sigma_{68}$ (MeV) & Core (\%) \\
\midrule
40 & 92,201 & 2 & +61.907 & 40.172 & 0.53 \\
60 & 243,668 & 6 & +7.500 & 40.204 & 39.18 \\
80 & 770,358 & 18 & -1.597 & 7.116 & 66.92 \\
100 & 1,309,730 & 35 & -2.411 & 3.724 & 72.61 \\
120 & 3,448,406 & 123 & -2.983 & 3.898 & 74.81 \\
140 & 1,966,199 & 82 & -3.510 & 4.441 & 75.64 \\
160 & 2,035,063 & 109 & -3.966 & 5.240 & 75.95 \\
180 & 1,998,339 & 129 & -4.382 & 6.311 & 75.86 \\
200 & 1,837,830 & 141 & -4.785 & 7.784 & 75.17 \\
220 & 1,619,560 & 206 & -5.250 & 9.868 & 73.84 \\
240 & 1,388,655 & 180 & -5.790 & 12.854 & 71.96 \\
260 & 1,157,665 & 176 & -6.581 & 17.255 & 69.90 \\
\bottomrule
\end{tabular}

\caption{Selected reconstructed MC entries and saved shape summaries.
Bias is the reconstructed median displacement, in MeV, from the generated
pole. Median and quantile half-width $\sigma_{68}$ are conditional on
the stored 0--400~MeV range; the core probability uses the full selected
weight. Overflow
denotes selected rows above that range; it remains included in
full-selected-yield bookkeeping.}
\end{table}

The original 2021 resolution is held fixed when defining every mask:
\begin{equation}
 \sigma_m(m)=0.00184825-0.001375m+0.085875m^2,
 \qquad m,\sigma_m\ \hbox{in GeV}.
\end{equation}
Thus a narrower or shifted MC core changes the fraction of a unit signal
inside the existing fitted bins. It does not silently redefine those bins
or the sideband training set. The MC mass range and normalization support
are recorded separately so losses at either boundary remain visible.

The word \emph{smeared} in the production directory is not sufficient
to reconstruct the upstream smearing chain. The stored smear-ratio
branches and cutflows are audited, while the template itself uses the
stored target-constrained mass exactly once. No second factor of 1.25
or undocumented mass transformation is introduced.

\clearpage
\section{Observed 2021 limits and local probabilities}
The direct MC points use the same observed counts and correlated GP
constraint as their Gaussian partners at each mass. The dense morph
curve is displayed with a separate line style because its additional
mass hypotheses require an interpolation assumption.

\fig{.73}{../figures/observed_2021.pdf}{Observed 2021 10\% limits,
fixed-mass local probabilities and signal fractions in the original
fitted bins. Points use the direct reconstructed templates at generated
masses. The dashed curve, when present, uses exploratory MC morphing.
The coupling limit includes the inherited minimal-visible branching
correction once above the dimuon threshold.}

\clearpage
\section{Paired comparison at generated masses}
The direct result is paired with the Gaussian at the same mass. The
limit ratio is $R(m)=\eps^2_{90,\mathrm{MC}}/\eps^2_{90,G}$, and the
signed-root change is $\Delta r=r_{\mathrm{MC}}-r_G$. Points at
different masses share observations and GP training bins.

\begin{table}[H]\centering\footnotesize
\begin{tabular}{rrrrrrr}
\toprule
$m_0$ & $\eps^2_{90,G}$ & $\eps^2_{90,MC}$ & Ratio & $p_{0,G}$ & $p_{0,MC}$ & $\Delta r$ \\
\midrule
60 & $3.747\times10^{-6}$ & $1.408\times10^{-5}$ & 3.7590 & $5.000\times10^{-1}$ & $4.950\times10^{-1}$ & +0.825 \\
80 & $5.221\times10^{-6}$ & $6.362\times10^{-6}$ & 1.2184 & $1.628\times10^{-2}$ & $2.477\times10^{-2}$ & -0.174 \\
100 & $2.599\times10^{-6}$ & $2.089\times10^{-6}$ & 0.8039 & $3.173\times10^{-1}$ & $5.000\times10^{-1}$ & -0.969 \\
120 & $2.905\times10^{-6}$ & $2.258\times10^{-6}$ & 0.7773 & $2.500\times10^{-1}$ & $5.000\times10^{-1}$ & -0.969 \\
140 & $1.968\times10^{-6}$ & $1.964\times10^{-6}$ & 0.9975 & $5.000\times10^{-1}$ & $5.000\times10^{-1}$ & -0.356 \\
160 & $5.376\times10^{-6}$ & $4.787\times10^{-6}$ & 0.8904 & $6.817\times10^{-2}$ & $2.188\times10^{-1}$ & -0.713 \\
180 & $5.946\times10^{-6}$ & $3.170\times10^{-6}$ & 0.5332 & $1.371\times10^{-1}$ & $5.000\times10^{-1}$ & -2.194 \\
200 & $5.577\times10^{-6}$ & $5.453\times10^{-6}$ & 0.9777 & $4.128\times10^{-1}$ & $5.000\times10^{-1}$ & -0.684 \\
220 & $6.114\times10^{-6}$ & $1.230\times10^{-5}$ & 2.0113 & $5.000\times10^{-1}$ & $5.000\times10^{-1}$ & +1.117 \\
240 & $1.558\times10^{-5}$ & $2.472\times10^{-5}$ & 1.5865 & $4.865\times10^{-1}$ & $4.436\times10^{-1}$ & +0.108 \\
\bottomrule
\end{tabular}

\caption{Direct-template comparison at generated masses. All
probabilities are pointwise and use their stated signal template.}
\end{table}

\fig{.39}{../figures/native_agreement.pdf}{Observed fractional limit
changes and signed-root differences for direct MC templates. These are
paired descriptions of the same observed spectra. They are not
independent residuals with a chi-squared sampling law.}

\clearpage
\section{Agreement summaries and interpolation boundary}
\begin{table}[H]\centering\small
\begin{tabular}{lr}
\toprule
Descriptive quantity & Direct MC versus Gaussian \\
\midrule
Common generated masses & 10 \\
Median ratio & 0.98760 \\
16th--84th percentile ratio & 0.78900--1.82440 \\
Minimum / maximum ratio & 0.53322 / 3.75896 \\
Median absolute relative change & 22.055\% \\
RMS log ratio & 0.54902 \\
Within 1 / 5 / 10\% & 10.0\% / 20.0\% / 20.0\% \\
Maximum $|\Delta r|$ & 2.19398 \\
\bottomrule
\end{tabular}

\caption{Descriptive agreement on the common direct-MC grid. Quantile
ranges and percentages of agreeing masses are summaries of the curves,
not confidence intervals or agreement probabilities.}
\end{table}

A statistical probability for disagreement would require a specified
joint sampling calculation that preserves the common data, overlapping
sidebands and template estimation. An independent-point chi-squared
test would omit those dependencies. The present comparison therefore
reports sizes and locations of differences without assigning such a
probability.

The reconstructed MC is available at discrete generated masses. A dense
scan necessarily interpolates shape information. Holding out an interior
generated mass and predicting its shape from its neighbors tests that
interpolation against an available simulation sample. It does not test
the underlying detector simulation against data, and it cannot establish
sampling coverage of an upper limit.

The declared morph linearly mixes the neighboring CDFs in standardized
residual $z=(m_{\rm rec}-m_0)/\sigma_{\rm nominal}(m_0)$, evaluating
each input distribution at its own pole and nominal resolution. It uses
no extrapolation beyond the supplied 40--260~MeV grid. Its dense observed
scan covers the original 50--250~MeV search at 1~MeV spacing.

8 of 10 held-out shape checks pass the declared $D_{\rm CDF}\leq0.03$ and $|\Delta F_{\rm core}|\leq0.02$ criteria. The largest conditional-support CDF distance is 0.1395. The held-out limit ratios to direct MC range from 0.9142 to 1.3029. The dense curve remains exploratory regardless of this outcome.


The direct generated-mass results remain the primary comparison. Any
minimum found on the denser exploratory grid depends on both the
morphing prescription and the enlarged mass scan. Neither its local
probability nor the smallest value selected across models is a global
discovery probability.

\clearpage
\subsection{Held-out shapes and finite-MC precision}
\begin{table}[H]\centering\small
\begin{tabular}{rrrrrl}
\toprule
$m_0$ & $D_{\rm CDF}$ & $\Delta F_{\rm core}$ & UL ratio & $\Delta r$ & Shape check \\
\midrule
60 & 0.1168 & -0.0545 & 0.9142 & +0.3724 & fail \\
80 & 0.1395 & -0.1090 & 1.3029 & +0.0490 & fail \\
100 & 0.0271 & -0.0166 & 1.0540 & +0.0364 & pass \\
120 & 0.0118 & -0.0062 & 1.0230 & +0.0116 & pass \\
140 & 0.0091 & -0.0025 & 1.0117 & +0.0037 & pass \\
160 & 0.0057 & -0.0016 & 1.0143 & +0.0103 & pass \\
180 & 0.0037 & -0.0028 & 1.0165 & +0.0124 & pass \\
200 & 0.0033 & -0.0030 & 1.0056 & -0.0154 & pass \\
220 & 0.0027 & -0.0026 & 1.0177 & +0.0289 & pass \\
240 & 0.0052 & -0.0006 & 1.0024 & +0.0168 & pass \\
\bottomrule
\end{tabular}

\caption{Leave-one-mass-out checks. CDF distances and core fractions
use the conditional analysis-support distributions. Limit ratios and
signed-root changes compare the held-out morph with direct MC at the
same hypothesis. The shape thresholds are diagnostic criteria.}
\end{table}

For each direct mass, 32 independent Poisson resamples of the stored
unit-weight MC histogram bins are fitted against the fixed Gaussian
reference. This conditional calculation quantifies template sampling
precision under that bin-count model. It does not include detector
systematics, selection differences, correlations between candidates or
uncertainty in the observed-data background. No new observed-data
experiments are generated.

\begin{table}[H]\centering\small
\begin{tabular}{rrrrrr}
\toprule
$m_0$ & Resamples & $R_{16}$ & $R_{50}$ & $R_{84}$ & UL spread (\%) \\
\midrule
60 & 32 & 3.7284 & 3.7609 & 3.7925 & 1.063 \\
80 & 32 & 1.2141 & 1.2173 & 1.2213 & 0.349 \\
100 & 32 & 0.8016 & 0.8037 & 0.8061 & 0.302 \\
120 & 32 & 0.7767 & 0.7772 & 0.7780 & 0.096 \\
140 & 32 & 0.9960 & 0.9978 & 0.9987 & 0.125 \\
160 & 32 & 0.8887 & 0.8908 & 0.8920 & 0.212 \\
180 & 32 & 0.5320 & 0.5332 & 0.5350 & 0.243 \\
200 & 32 & 0.9750 & 0.9773 & 0.9798 & 0.268 \\
220 & 32 & 2.0067 & 2.0148 & 2.0185 & 0.312 \\
240 & 32 & 1.5794 & 1.5852 & 1.5907 & 0.326 \\
\bottomrule
\end{tabular}

\caption{Conditional finite-MC precision. $R_{16},R_{50},R_{84}$ are
quantiles of the MC-to-Gaussian upper-limit ratio across 32 MC-bin
resamples. The last column is the relative standard deviation of the
resampled limit around its central direct-template scale, in percent.}
\end{table}

Thirty-two resamples resolve only a modest estimate of the central
spread. The intervals are neither an observed-limit confidence band nor
a significance for the discrepancy between templates.

\clearpage
\section{Likelihood and normalization}
The baseline uses the original primary window. The GP trains only on
bins outside $\pm2.25\sigma_m$, using archived kernel coordinates.
The window diagnostics replace 2.25 by 2.4, 2.5 or 2.6 for both the fit
and training mask, recomputing the prediction with the same kernel states.
Within a mass hypothesis, both signal models share its predicted count
mean $b$, covariance $C$ and observed counts $n$. With $LL^T=C$,
\begin{align}
 a&=\frac{\eps^2}{10^{-8}},\qquad T_i=K(m)\,10^{-8}w_i,\\
 \lambda_i&=b_i+(L\bm\theta)_i+a T_i,\\
 \ell(a,\bm\theta)&=\sum_i\left[\lambda_i-n_i+
          n_i\log\frac{n_i}{\lambda_i}\right]
          +\frac12\bm\theta^T\bm\theta .
\end{align}
The zero-count contribution is $\lambda_i$. The correlated background
deformation is profiled at every tested amplitude. The bin probabilities $w_i$ are
normalized over the declared full support before selecting fitted bins;
their restriction to the primary window is not renormalized.

For MC the full fitting support is 36--299.75~MeV. Its retained fraction
$f_{\rm support}$ is recorded separately. The ledger's
\texttt{full\_yield\_90} is the support-normalized yield; the additional
\texttt{full\_selected\_yield\_90} divides it by $f_{\rm support}$.
That bookkeeping conversion does not establish a calibrated coupling
normalization outside the analysis support.

Here $K(m)$ is the unchanged prompt-density conversion. The solver's
\texttt{A90} column is $a_{90}$, while the support-normalized event yield
is $N_{90}=a_{90}K(m)10^{-8}=K(m)\eps^2_{90}$. Above $2m_\mu$, the inherited
minimal-visible conversion multiplies the electron-only coupling limit
by $1+\sqrt{1-4r_\mu}(1+2r_\mu)$, with
$r_\mu=(m_\mu/m)^2$. This correction does not change the visible yield
or local probability. Density continues to use the nominal
$m\pm1.64\sigma_m$ native-bin interval.

The observed limit solves the bounded, piecewise-asymptotic
$CL_s(a_{90})=0.10$, using explicitly profiled observed and
background-Asimov likelihood ratios~\cite{Cowan}. The free diagnostic
amplitude may be negative while physical Poisson expectations are
maintained. Writing its signed likelihood root as $r$,
\begin{equation}
 Z_0=\max(r,0),\qquad p_0=\Phi(-Z_0).
\end{equation}
Deficits therefore have the inherited reference value $p_0=0.5$.
This is a local asymptotic reference, not the probability of the most
significant result anywhere in a scan.

The Gaussian and MC comparisons use the same converged solver. The
historical v5.9 audit found a stopping discrepancy in an older optimizer;
the present paired comparison does not mix that older table with newly
computed MC fits. Numerical convergence and finite-MC statistical
qualification are separate questions.

\clearpage
\section{Selection provenance and interpretation}
The supplied MC directory is
\path{/sdf/data/hps/physics2021/preselection/v13/ap_signal_prompt_smeared}.
The observed input derives from
\path{/sdf/data/hps/physics2021/preselection/v16/data_10pc_prompt_TC_psum2p8}.
The MC uses the stored \texttt{vertex./vertex.invM\_} branch and the
target-constrained vertex category. Production labels agree on the
target constraint and $p_{\rm sum}>2.8$~GeV, but the inspected MC timing
windows include 5.7 and 9.2~ns where the observed sample uses 5.1 and
7.8~ns. Full production equivalence is not established. No new selection
is silently imposed to make the two samples appear matched.

A bounded inspection finds both stored daughter truth-link flags false
for every inspected row. Those flags cannot identify matched signal
daughters, and their values alone do not prove that the reconstructed
candidates lack signal daughters. No additional truth selection is
defensible from them. The fit therefore substitutes the supplied
all-selected candidate distribution conditionally.

The stored track hit counters are three-dimensional; cutflow labels
refer to two-dimensional-hit requirements. Those labels are not a
license to apply their thresholds directly to the stored hit counters.
The local and remote selection audits preserve the evidence and the
unresolved production details. ROOT files can contain older tree cycles;
only the latest \texttt{preselection} cycle contributes to extraction.

Input files and configurations are identified by saved paths, sizes and
SHA-256 hashes. Extraction uses bounded chunks and retains bin sums of
weights and squared weights. All results, scripts, reconstructed-mass
histograms and report sources are mirrored to
\path{/sdf/home/e/epeets/src/hps_gpr_v6p1_mc_signal_templates_20260922}.
The original observed result and source ROOT files are preserved.

This study measures how the selected reconstructed 2021 MC changes the
conditional observed inference under the archived background model and
the explicitly declared fit and training mask. It does not isolate a pure tail uncertainty or establish that
either signal template is the correct detector response. Selection
differences, finite MC and interpolation remain explicit qualifications.
No MC-template claim is made for the 2015 or 2016 detector configurations.

Approximately 24--61\% of the direct-mass selected signal distributions
lie outside the primary blind window. Signal in those training bins can
be absorbed by a background-only GP. This observed template replacement
keeps the sideband prediction fixed; it does not model signal-contamination
feedback, extraction recovery or coverage for these broad distributions.

\begin{thebibliography}{9}
\bibitem{note} E.~Peets, M.~Gignac and E.~Hsu,
\emph{HPS Gaussian-Process Resonance Search: Analysis Note v5.0.5},
16 September 2026, internal review draft and archived source package.
\bibitem{Cowan} G.~Cowan, K.~Cranmer, E.~Gross and O.~Vitells,
``Asymptotic formulae for likelihood-based tests of new physics,''
Eur. Phys. J. C \textbf{71}, 1554 (2011),
\href{https://arxiv.org/abs/1007.1727}{arXiv:1007.1727}.
\end{thebibliography}
\clearpage
\section{Common fit and training windows of 2.4, 2.5 and 2.6 sigma}
The original $\pm2.25\sigma_m$ prescription remains the baseline. Three
additional calculations use common fit and training-mask half-widths
$h=2.4\sigma_m$, $2.5\sigma_m$ and $2.6\sigma_m$. There is no separate
outer guard window: the same half-width determines the fitted bins and
the excluded GP training bins. Each changed mask recomputes the GP
prediction and correlated covariance from its own observed sidebands,
with the archived kernel coordinates held fixed.

Gaussian and MC templates at a given width share those counts and GP
constraints. A comparison between widths therefore changes both the
observed fit region and the background prediction; it cannot isolate
signal-tail acceptance alone. The nominal resolution, coupling conversion
and full-support template normalization are retained. Masks select bins by their native-bin centers. Nearby half-widths
can therefore select exactly the same fit and training bins at a particular
mass and yield coincident results; the width response is not continuous. No width is selected for
its observed limit or local probability.

\begin{table}[H]\centering\scriptsize
\begin{tabular}{rrrrrrr}
\toprule
$h/\sigma_m$ & Median ratio & Ratio range & Med. $|R-1|$ (\%) & Max. $|\Delta r|$ & LOO pass & MC spread (\%) \\
\midrule
2.25 & 0.9876 & 0.5332--3.7590 & 22.05 & 2.194 & 8/10 & 1.06 \\
2.4 & 1.0038 & 0.5413--3.7386 & 22.92 & 2.164 & 8/10 & 1.11 \\
2.5 & 1.0127 & 0.5480--3.9413 & 28.90 & 2.151 & 8/10 & 1.13 \\
2.6 & 1.0127 & 0.5491--3.9413 & 29.53 & 2.170 & 8/10 & 1.13 \\
\bottomrule
\end{tabular}

\caption{Direct-MC versus Gaussian comparisons at ten generated masses,
reported separately for each common half-width. The LOO column counts
held-out shape checks satisfying the declared criteria. MC spread is
the largest relative upper-limit standard deviation among the ten
masses under 32 independent-bin Poisson resamples. These descriptions
are correlated across both mass and window.}
\end{table}

The median absolute MC-to-Gaussian limit change is 22.92\%, 28.90\%
and 29.53\% for half-widths 2.4, 2.5 and 2.6 sigma, compared with
22.05\% at baseline. The smallest direct-template local probability
remains at 80~MeV: $p_0=0.02477$ for baseline and 2.4 sigma, and
$p_0=0.009903$ for 2.5 and 2.6 sigma. These are conditional local
results, not a basis for choosing the width.

The following pages keep the three new prescriptions separate. Their
upper panels show the full observed curves; the lower panels and tables
compare direct templates only at their generated masses. Dashed curves
use the same exploratory CDF morph as the baseline, including its
low-mass qualification. Local probabilities retain $p_0=0.5$ for
deficits and are not adjusted for inspecting multiple masses, models
or window widths.

The finite-MC and leave-one-out diagnostics are recomputed for each
width and retained in the portable ledgers. Wider masks do not resolve
the unverified candidate association, production-selection mismatch or
signal leakage into the sidebands. The inherited coupling values remain
conditional template comparisons rather than validated A-prime exclusions.

\clearpage
\subsection{Half-width $2.4\sigma_m$}
\fig{.55}{../figures/window_2p4.pdf}{2021 10\% observed results with
common fit and GP training-mask half-width $2.4\sigma_m$. The Gaussian
and MC predictions share the background constraint within this page.
The lower panels show direct-MC/Gaussian limit ratios and signed-root
changes at generated masses; dashed dense curves are exploratory.}
\begin{table}[H]\centering\scriptsize
\begin{tabular}{rrrrrrr}
\toprule
$m_0$ & $\eps^2_{90,G}$ & $\eps^2_{90,MC}$ & Ratio & $p_{0,G}$ & $p_{0,MC}$ & $\Delta r$ \\
\midrule
60 & $3.73\times10^{-6}$ & $1.39\times10^{-5}$ & 3.739 & 0.5000 & 0.5000 & +0.616 \\
80 & $5.22\times10^{-6}$ & $6.36\times10^{-6}$ & 1.218 & 0.0163 & 0.0248 & -0.174 \\
100 & $2.60\times10^{-6}$ & $2.09\times10^{-6}$ & 0.804 & 0.3173 & 0.5000 & -0.969 \\
120 & $2.95\times10^{-6}$ & $2.24\times10^{-6}$ & 0.760 & 0.2465 & 0.5000 & -1.053 \\
140 & $1.94\times10^{-6}$ & $1.95\times10^{-6}$ & 1.007 & 0.5000 & 0.5000 & -0.318 \\
160 & $5.38\times10^{-6}$ & $4.56\times10^{-6}$ & 0.848 & 0.0716 & 0.2724 & -0.858 \\
180 & $6.23\times10^{-6}$ & $3.37\times10^{-6}$ & 0.541 & 0.1148 & 0.5000 & -2.164 \\
200 & $5.62\times10^{-6}$ & $5.62\times10^{-6}$ & 1.001 & 0.4175 & 0.5000 & -0.635 \\
220 & $6.25\times10^{-6}$ & $1.34\times10^{-5}$ & 2.141 & 0.5000 & 0.4879 & +1.256 \\
240 & $1.58\times10^{-5}$ & $2.48\times10^{-5}$ & 1.569 & 0.4771 & 0.4595 & +0.044 \\
\bottomrule
\end{tabular}

\caption{Direct-template comparison for half-width $2.4\sigma_m$.
Coupling limits and local probabilities retain the same conventions
as the baseline table.}
\end{table}

\clearpage
\subsection{Half-width $2.5\sigma_m$}
\fig{.55}{../figures/window_2p5.pdf}{2021 10\% observed results with
common fit and GP training-mask half-width $2.5\sigma_m$. The Gaussian
and MC predictions share the background constraint within this page.
The lower panels show direct-MC/Gaussian limit ratios and signed-root
changes at generated masses; dashed dense curves are exploratory.}
\begin{table}[H]\centering\scriptsize
\begin{tabular}{rrrrrrr}
\toprule
$m_0$ & $\eps^2_{90,G}$ & $\eps^2_{90,MC}$ & Ratio & $p_{0,G}$ & $p_{0,MC}$ & $\Delta r$ \\
\midrule
60 & $3.80\times10^{-6}$ & $1.50\times10^{-5}$ & 3.941 & 0.5000 & 0.5000 & +0.727 \\
80 & $5.47\times10^{-6}$ & $7.31\times10^{-6}$ & 1.338 & 0.0121 & 0.0099 & +0.078 \\
100 & $2.65\times10^{-6}$ & $2.20\times10^{-6}$ & 0.831 & 0.3101 & 0.5000 & -0.915 \\
120 & $2.95\times10^{-6}$ & $2.24\times10^{-6}$ & 0.760 & 0.2465 & 0.5000 & -1.053 \\
140 & $1.94\times10^{-6}$ & $2.03\times10^{-6}$ & 1.047 & 0.5000 & 0.5000 & -0.225 \\
160 & $5.44\times10^{-6}$ & $4.62\times10^{-6}$ & 0.849 & 0.0698 & 0.2743 & -0.878 \\
180 & $6.24\times10^{-6}$ & $3.42\times10^{-6}$ & 0.548 & 0.1170 & 0.5000 & -2.151 \\
200 & $5.60\times10^{-6}$ & $5.47\times10^{-6}$ & 0.978 & 0.4234 & 0.5000 & -0.725 \\
220 & $6.18\times10^{-6}$ & $1.29\times10^{-5}$ & 2.084 & 0.5000 & 0.5000 & +1.163 \\
240 & $1.62\times10^{-5}$ & $2.74\times10^{-5}$ & 1.686 & 0.4545 & 0.3763 & +0.201 \\
\bottomrule
\end{tabular}

\caption{Direct-template comparison for half-width $2.5\sigma_m$.
Coupling limits and local probabilities retain the same conventions
as the baseline table.}
\end{table}

\clearpage
\subsection{Half-width $2.6\sigma_m$}
\fig{.55}{../figures/window_2p6.pdf}{2021 10\% observed results with
common fit and GP training-mask half-width $2.6\sigma_m$. The Gaussian
and MC predictions share the background constraint within this page.
The lower panels show direct-MC/Gaussian limit ratios and signed-root
changes at generated masses; dashed dense curves are exploratory.}
\begin{table}[H]\centering\scriptsize
\begin{tabular}{rrrrrrr}
\toprule
$m_0$ & $\eps^2_{90,G}$ & $\eps^2_{90,MC}$ & Ratio & $p_{0,G}$ & $p_{0,MC}$ & $\Delta r$ \\
\midrule
60 & $3.80\times10^{-6}$ & $1.50\times10^{-5}$ & 3.941 & 0.5000 & 0.5000 & +0.727 \\
80 & $5.47\times10^{-6}$ & $7.31\times10^{-6}$ & 1.338 & 0.0121 & 0.0099 & +0.078 \\
100 & $2.65\times10^{-6}$ & $2.20\times10^{-6}$ & 0.831 & 0.3101 & 0.5000 & -0.915 \\
120 & $2.97\times10^{-6}$ & $2.22\times10^{-6}$ & 0.747 & 0.2454 & 0.5000 & -1.131 \\
140 & $1.94\times10^{-6}$ & $2.03\times10^{-6}$ & 1.047 & 0.5000 & 0.5000 & -0.225 \\
160 & $5.47\times10^{-6}$ & $4.78\times10^{-6}$ & 0.873 & 0.0684 & 0.2549 & -0.829 \\
180 & $6.20\times10^{-6}$ & $3.40\times10^{-6}$ & 0.549 & 0.1234 & 0.5000 & -2.170 \\
200 & $5.65\times10^{-6}$ & $5.52\times10^{-6}$ & 0.978 & 0.4173 & 0.5000 & -0.723 \\
220 & $6.16\times10^{-6}$ & $1.29\times10^{-5}$ & 2.090 & 0.5000 & 0.5000 & +1.146 \\
240 & $1.62\times10^{-5}$ & $2.71\times10^{-5}$ & 1.670 & 0.4606 & 0.4053 & +0.141 \\
\bottomrule
\end{tabular}

\caption{Direct-template comparison for half-width $2.6\sigma_m$.
Coupling limits and local probabilities retain the same conventions
as the baseline table.}
\end{table}

\clearpage
\section{Morphed midpoint templates: full reconstructed range}
The eleven midpoint hypotheses use equal CDF mixtures of their two
neighboring generated samples, transported through the nominal-resolution
standardized coordinate. The plots use saved numerical morphs, not a fit
to observed data. The bracketing curves show the transported inputs at
the midpoint hypothesis. Each displayed curve is normalized over
0--400~MeV; fitted probabilities retain the separately recorded analysis
support normalization.

\fig{.74}{../figures/morph_gallery_full.pdf}{Full-range reconstructed
mass densities for midpoint morphs at 50, 70, $\ldots$, 250~MeV, with
transported neighboring MC shapes. The logarithmic scale exposes broad
and displaced populations. Dotted lines mark the baseline
$\pm2.25\sigma_m$ interval for reference; the morph itself is independent
of the fit-window width. Asterisks flag the 50/70/90~MeV hypotheses that
touch the low-mass region with failed held-out shape checks.}

\clearpage
\section{Morphed midpoint templates: core detail}
The same saved distributions are displayed within seven nominal sigma
of each midpoint, on a linear density scale. Zooming the horizontal
axis does not renormalize the curves or remove their broad components
from the fitted signal model. The interpolation supplies no additional
truth matching, production-equivalence validation or detector-response
certification.

\fig{.76}{../figures/morph_gallery_core.pdf}{Core detail of the same
midpoint morphs and transported neighboring shapes. The plot scale is
set independently in each panel. Dotted boundaries retain the original
baseline half-width. These MC-only interpolations remain exploratory,
particularly the marked low-mass hypotheses; they are not new simulated
samples.}

\clearpage
\section{MC core centers and shifted primary windows}
\label{sec:corecenter}
This update, dated 23 September 2026, centers the observed fit and the
excluded GP training interval on the reconstructed MC core, rather than
the generated pole. The reconstructed template itself is unchanged.
For a generated hypothesis $m_0$ and an MC-only core location $c(m_0)$,
\begin{equation}
 I(m_0)=[c(m_0)-2.25\sigma_{\rm nominal}(m_0),\,
         c(m_0)+2.25\sigma_{\rm nominal}(m_0)].
\end{equation}
The half-width retains the original nominal resolution. The fitted core
width below does not replace it. No additional guard is used.

The center is determined without observed counts. A Gaussian smoothing
of width $0.2\sigma_{\rm nominal}$ locates the largest MC peak within
$m_0\pm3\sigma_{\rm nominal}$. A bin-integrated Gaussian plus a
nonnegative affine pedestal is then fitted by Poisson deviance to the
unsmoothed template within $\pm1.5\sigma_{\rm nominal}$ of that mode.
The Gaussian mean defines $c$; the pedestal limits the influence of a
local broad component. The mean may vary by one nominal sigma around
the mode, and the fitted width by 0.25--2.5 nominal sigma. All 201 fits
converge without these bounds being active. This is an operational core
definition, not proof of a pure resonance response.

\begin{table}[H]\centering\scriptsize
\begin{tabular}{rrrrrrr}
\toprule
$m_0$ & $c(m_0)$ & $c-m_0$ & $\sigma_{\rm core}$ & MC spread & Max. $|\Delta c|$ & Max. UL change (\%) \\
\midrule
60 & 58.792 & -1.208 & 1.550 & 0.034 & 0.049 & 0.00 \\
80 & 78.103 & -1.897 & 1.637 & 0.007 & 0.003 & 0.00 \\
100 & 97.688 & -2.312 & 1.853 & 0.006 & 0.027 & 0.00 \\
120 & 117.319 & -2.681 & 2.112 & 0.003 & 0.039 & 0.00 \\
140 & 136.956 & -3.044 & 2.392 & 0.006 & 0.052 & 0.29 \\
160 & 156.658 & -3.342 & 2.713 & 0.007 & 0.081 & 0.00 \\
180 & 176.393 & -3.607 & 3.120 & 0.008 & 0.096 & 0.00 \\
200 & 196.182 & -3.818 & 3.577 & 0.011 & 0.130 & 0.00 \\
220 & 215.941 & -4.059 & 4.233 & 0.012 & 0.162 & 4.01 \\
240 & 235.617 & -4.383 & 4.981 & 0.022 & 0.132 & 3.75 \\
\bottomrule
\end{tabular}

\caption{Native-MC core centers and stability. All mass quantities are
in MeV. MC spread is the standard deviation of the estimated center
under 32 independent-bin Poisson resamples. The last two columns vary
the core-fit half-width from its default 1.5 nominal sigma to 1.25 and
2.0, and refit the observed MC limit using each alternative center.}
\end{table}

The maximum center-definition displacement is 0.162~MeV and the largest
corresponding limit change is 4.01\%. Native-bin masks are discrete;
many smaller shifts select exactly the same bins. Bin resampling
quantifies conditional center precision only; it does not establish
candidate independence or cover production systematics.

The direct-template centers shift by 1.21--4.38~MeV below their generated poles. The median absolute MC/Gaussian limit difference in the matched centered windows is 17.83\%, compared with 22.05\% for the original paired prescription. The native new/old MC limit ratios span 0.822--1.201. The smallest native MC local probability is at $m_0=80$~MeV, $c=78.103$~MeV: $p_0=0.00787$ ($Z_0=2.415$), compared with $p_0=0.02477$ before shifting the window.


\clearpage
\subsection{Observed scans after centering}
The GP mean and covariance are recomputed from sidebands outside each
shifted interval. Archived kernel coordinates remain anchored at $m_0$.
The inherited density conversion, nominal resolution and branching
conversion also retain $m_0$, isolating the center change from a change
of coupling convention. Both the MC and its new Gaussian benchmark use
the same shifted window and GP constraint. The benchmark Gaussian is
centered at $c(m_0)$ with width $\sigma_{\rm nominal}(m_0)$; it is not
the historical Gaussian fit evaluated at mass $c$ with new kernel or
density settings. All curves are plotted against generated hypothesis
$m_0$, not reconstructed center $c$.

\fig{.71}{../figures/core_centered_comparison.pdf}{Original and
MC-centered observed inference for the 2021 10\% sample. Native MC
points are marked; continuous MC curves include exploratory morphs.
The bottom panels compare only the ten generated masses. The matched
Gaussian comparison separates the remaining shape difference from
placing a Gaussian peak at the wrong reconstructed position.}

\clearpage
\subsection{Native-mass limits and local probabilities}
\begin{table}[H]\centering\scriptsize
\begin{tabular}{rrrrrrr}
\toprule
$m_0$ & Old MC limit & New MC limit & New/old & Old MC $p_0$ & New MC $p_0$ & New G $p_0$ \\
\midrule
60 & $1.41\times10^{-5}$ & $1.16\times10^{-5}$ & 0.822 & 0.49502 & 0.50000 & 0.50000 \\
80 & $6.36\times10^{-6}$ & $7.64\times10^{-6}$ & 1.201 & 0.02477 & 0.00787 & 0.00348 \\
100 & $2.09\times10^{-6}$ & $2.30\times10^{-6}$ & 1.103 & 0.50000 & 0.50000 & 0.50000 \\
120 & $2.26\times10^{-6}$ & $2.30\times10^{-6}$ & 1.021 & 0.50000 & 0.50000 & 0.50000 \\
140 & $1.96\times10^{-6}$ & $2.09\times10^{-6}$ & 1.067 & 0.50000 & 0.50000 & 0.50000 \\
160 & $4.79\times10^{-6}$ & $5.04\times10^{-6}$ & 1.053 & 0.21877 & 0.21609 & 0.17038 \\
180 & $3.17\times10^{-6}$ & $3.26\times10^{-6}$ & 1.028 & 0.50000 & 0.50000 & 0.50000 \\
200 & $5.45\times10^{-6}$ & $5.22\times10^{-6}$ & 0.957 & 0.50000 & 0.50000 & 0.50000 \\
220 & $1.23\times10^{-5}$ & $1.21\times10^{-5}$ & 0.981 & 0.50000 & 0.50000 & 0.50000 \\
240 & $2.47\times10^{-5}$ & $2.32\times10^{-5}$ & 0.940 & 0.44359 & 0.50000 & 0.50000 \\
\bottomrule
\end{tabular}

\caption{Original and shifted-window MC results. Limits use the inherited
coupling coordinate. The last column is the Gaussian benchmark centered
at the same MC core. All probabilities are local asymptotic references.}
\end{table}
\begin{table}[H]\centering\small
\begin{tabular}{rrrrrr}
\toprule
$m_0$ & Gaussian at $c$ limit & MC/G at $c$ & $r_{MC}-r_G$ & Old MC fit (\%) & New MC fit (\%) \\
\midrule
60 & $4.54\times10^{-6}$ & 2.548 & -0.520 & 39.26 & 39.65 \\
80 & $6.20\times10^{-6}$ & 1.232 & -0.284 & 67.56 & 70.95 \\
100 & $2.19\times10^{-6}$ & 1.052 & -0.292 & 73.19 & 78.21 \\
120 & $2.12\times10^{-6}$ & 1.088 & -0.076 & 74.60 & 80.37 \\
140 & $1.84\times10^{-6}$ & 1.136 & +0.068 & 76.52 & 81.30 \\
160 & $4.76\times10^{-6}$ & 1.059 & -0.167 & 75.22 & 81.75 \\
180 & $3.27\times10^{-6}$ & 0.997 & -0.434 & 75.27 & 80.46 \\
200 & $4.27\times10^{-6}$ & 1.221 & -0.002 & 75.71 & 79.37 \\
220 & $8.91\times10^{-6}$ & 1.355 & +0.073 & 73.86 & 77.30 \\
240 & $1.63\times10^{-5}$ & 1.427 & -0.048 & 72.38 & 74.97 \\
\bottomrule
\end{tabular}

\caption{Matched centered-window comparison and retained MC probability
in the fitted bins. Support normalization remains unchanged; shifting
the mask does not renormalize its retained template fraction to one.}
\end{table}

At the native 80~MeV hypothesis, the reconstructed center is 78.10~MeV.
The shifted MC upper limit is $7.639\times10^{-6}$, compared with
$6.362\times10^{-6}$ before shifting, a 20.08\% increase. The local
probability decreases from 0.02477 to 0.007875; an increased upper limit
and a smaller local probability are compatible because they answer
different likelihood questions.

On the dense exploratory MC grid the smallest local probability moves
to 50~MeV ($p_0=0.003798$). That template uses the previously failing
low-mass interpolation region. It is not promoted to a validated
search result; the native comparison above remains the primary result.
The unchanged selected-candidate association and v13/v16 production
qualifications still apply. No global probability, coverage calibration
or physically validated A-prime exclusion follows from this recentering.

\clearpage
\subsection{Reconstructed cores and window placement}
\fig{.76}{../figures/core_centered_windows.pdf}{Native selected-MC cores
and the original versus shifted $\pm2.25\sigma_{\rm nominal}$ windows.
Blue curves are unchanged empirical templates. The local Gaussian plus
affine-pedestal fit (red) estimates the center only and is never used to
replace the full MC signal shape in the observed likelihood. The red
vertical line marks the fitted core center.}

\clearpage
\section{An analytic description of the MC core shift}
The ten native cores within the 60--240~MeV study range support a compact empirical
location law. Four simple forms are fitted with equal weight per mass:
a proportional shift, an affine function, a quadratic and a logarithm.
Each mass is also omitted in turn and predicted from the other nine.
No observed spectrum enters this comparison. The best leave-one-mass-out
(LOO) RMS among these candidates is the logarithmic model,
\begin{equation}
 c(m)-m=-3.22433-2.21399\ln\!\left(\frac{m}{150}\right)\ \mathrm{MeV},
 \qquad 60\leq m\leq240~\mathrm{MeV}.
\end{equation}
The argument uses $m$ in MeV. The law is an empirical approximation,
not a detector mechanism or an extrapolation outside the fitted range.
Its fitted RMS is 0.0497~MeV, LOO RMS 0.0652~MeV and largest held-out
error 0.1604~MeV. These errors describe cross-mass prediction, not a
calibrated confidence band. Model selection and evaluation use the same
ten masses; the LOO summary is a diagnostic, not independent validation.
\fig{.32}{../figures/analytic_center_models.pdf}{Measured core shifts,
analytic models and their held-out residuals. Small error bars show
conditional MC-bin resampling spread; the equal-mass fits do not treat
that spread as the complete response uncertainty.}
\begin{table}[H]\centering\small
\begin{tabular}{lrrrr}
\toprule
Model & Parameters & Fit RMS (MeV) & LOO RMS (MeV) & Max. LOO (MeV) \\
\midrule
proportional & 1 & 0.2552 & 0.2885 & 0.4602 \\
affine & 2 & 0.1555 & 0.2154 & 0.5287 \\
quadratic & 3 & 0.0679 & 0.1399 & 0.3015 \\
logarithmic & 2 & 0.0497 & 0.0652 & 0.1604 \\
\bottomrule
\end{tabular}

\caption{Analytic model comparison. Affine and quadratic forms use
$x=(m-150)/100$; coefficients and all held-out predictions are saved.}
\end{table}

\clearpage
\section{Can one normalized core describe all masses?}
The ten distributions are first recentered without rescaling their
physical widths. They are then expressed in
$u=(m_{ee}-c)/\sigma_{\rm core}$ using the measured local Gaussian width.
The density transformation includes the Jacobian:
$q_m(u)=\sigma_{\rm core}p_m(c+\sigma_{\rm core}u)$.
Each full source histogram has unit area over its recorded 0--400~MeV
range. Display cropping never renormalizes the full distributions.
The separately labeled core view conditions on $|u|<2$ and therefore
discards information about the fraction of all candidates in that core.
\fig{.64}{../figures/shared_shape_overlays.pdf}{All ten normalized native
MC histograms superimposed; colored curves are empirical histograms,
not Gaussian fits or morphed samples. Top left: centered physical-mass densities. Top right:
full-normalized densities after width scaling. Bottom left: unit-area
conditional cores. Bottom right: full-normalized tails on a log scale.
The black dashed curve is the equal-mass mean standardized distribution;
sample event counts do not determine its weights. Direct MC-point
overlays and the rebinning definition appear in Section~\ref{sec:MCpoints}.}
The conditional cores are much closer than the full shapes. A shared
central shape is therefore a useful possibility, but its unit-area
overlay cannot establish a common core fraction or common tails. The
full candidate distributions include the same unresolved response and
candidate-association contributions as the earlier sections.

\clearpage
\subsection{Held-out tests of a common shape}
For each generated mass, the common standardized CDF is averaged from
the other nine histograms. The primary shape-only test supplies the held-out
sample's known center and width. A second test predicts its center from
the selected analytic law fitted to the other nine and predicts its
width from a quadratic in $\ln(\sigma_{\rm core}/\mathrm{MeV})$ versus
$x=(m-150)/100$, also fitted without that sample. All likelihood comparisons
use the same measured-center evaluation window and GP constraint as the
direct MC. The second test predicts the complete template but does not
also move its evaluation mask. Existing neighboring-mass morphs are
evaluated in that same window as a comparison.
\begin{table}[H]\centering\scriptsize
\begin{tabular}{rrrrrrr}
\toprule
$m_0$ & Core (\%) & $D_{core}$ & $D_{full}$ & Shared UL ratio & Morph UL ratio & Predictive UL ratio \\
\midrule
60 & 32.4 & 0.045 & 0.461 & 0.443 & 1.047 & 0.419 \\
80 & 60.7 & 0.009 & 0.075 & 0.944 & 1.324 & 0.991 \\
100 & 68.1 & 0.008 & 0.023 & 1.140 & 1.050 & 1.155 \\
120 & 71.0 & 0.010 & 0.061 & 1.191 & 1.025 & 1.172 \\
140 & 71.5 & 0.010 & 0.072 & 1.135 & 1.009 & 1.124 \\
160 & 70.7 & 0.009 & 0.067 & 1.135 & 1.016 & 1.127 \\
180 & 69.6 & 0.007 & 0.052 & 1.085 & 1.017 & 1.076 \\
200 & 67.5 & 0.007 & 0.033 & 1.002 & 1.003 & 1.003 \\
220 & 65.8 & 0.009 & 0.037 & 0.936 & 1.018 & 0.936 \\
240 & 63.6 & 0.014 & 0.049 & 0.864 & 1.005 & 0.860 \\
\bottomrule
\end{tabular}

\caption{Held-out CDF distances and observed upper-limit ratios relative
to direct MC. Core probability is conditional on the stored histogram
range; $D_{core}$ conditions on $|u|<2$. Full distances use the declared
analysis support. ``Predictive'' estimates center and width as well as
shape without the held-out histogram.}
\end{table}
\fig{.26}{../figures/shared_shape_validation.pdf}{Common-shape and
neighbor-morph prediction errors, with matched likelihood conditions.}
Core-conditioned CDF distances span 0.0067--0.0450, while full-support
distances span 0.0229--0.4611. The common-shape limit ratios span
0.443--1.191 even when location and width are supplied. Thus a single
full standardized shape is not supported as a drop-in replacement.
A shared core with mass-dependent core fraction and tails is better
motivated by these diagnostics. For inspection, 181 exploratory common-shape
templates are saved at 1~MeV spacing over 60--240~MeV using the analytic
center and width laws; they are not adopted as the nominal signal model.

\clearpage
\section{Left and right signal probability and asymmetric blinding}
Asymmetry is measured around the fitted core, retaining the full
stored-histogram normalization. The upper panel measures probability
beyond $\pm2.25\sigma_{\rm nominal}$; the lower panel compares the two
halves inside that interval. Source overflow is below 0.016\% and is
excluded by the explicitly declared 0--400~MeV normalization here.
\fig{.43}{../figures/left_right_MC_probability.pdf}{Mass dependence of
the left and right selected-MC contributions around the core.}
The left external tail is larger at 120--240~MeV, whereas the right
external tail is larger at 60--100~MeV. The direction is therefore not
universal across the accepted samples.

Five windows are declared around the same measured core before fitting
the observed data: $[-2.25,+2.25]$, $[-2.5,+2.0]$, $[-2.0,+2.5]$,
$[-2.75,+2.25]$ and $[-2.25,+2.75]$, all in nominal sigma.
The redistributed pair keeps total width 4.5 sigma; each is also exactly
a shift of the window midpoint by $\mp0.25$ sigma, so it cannot isolate
tail asymmetry from center placement. The expanded pair adds half a
sigma on either side as matched controls.

Each window defines both the fitted bins and excluded GP training bins.
The GP prediction and covariance are recomputed; archived kernel,
resolution and coupling conventions remain anchored at the generated
mass. Gaussian and direct-MC templates share the same mask and background.
Actual bin counts are saved because equal continuous widths need not
select the same number of discrete observed bins.

\clearpage
\subsection{Observed asymmetric-window comparisons}
\begin{table}[H]\centering\scriptsize
\begin{tabular}{lrrrr}
\toprule
Window in nominal sigma & MC UL / symmetric & Med. $|MC/G-1|$ (\%) & Min. mass & Min. local p0 \\
\midrule
$[-2.25,+2.25]$ & 1.000--1.000 & 17.8 & 80 & 0.00787 \\
$[-2.5,+2]$ & 0.942--1.038 & 19.4 & 80 & 0.00671 \\
$[-2,+2.5]$ & 0.964--1.119 & 18.6 & 80 & 0.00854 \\
$[-2.75,+2.25]$ & 0.983--1.129 & 20.5 & 80 & 0.00761 \\
$[-2.25,+2.75]$ & 0.997--1.133 & 18.5 & 80 & 0.00868 \\
\bottomrule
\end{tabular}

\caption{Prespecified native-mass window comparisons. The minimum local
probability is descriptive and is not corrected for testing several
windows. No prescription is chosen from the observed result.}
\end{table}
\fig{.48}{../figures/asymmetric_comparisons.pdf}{Native MC limit ratios,
signed-root changes, local probabilities and fitted signal fractions for
all five windows. Full MC/Gaussian limits and probabilities are retained
in the accompanying 100-row fit ledger. Changes include altered data bins
and altered background interpolation, as well as signal acceptance.}
These are conditional observed diagnostics. A preferred production
window would require signal-recovery and coverage checks, including
signal leakage into the training sidebands, rather than selection by
the largest observed local significance or tightest upper limit.

\clearpage
\section{What the Figure 12 curves represent}\label{sec:MCpoints}
The colored curves are the ten native reconstructed-MC histograms,
normalized and plotted as connected bin densities. The local Gaussian
plus affine-pedestal fit supplies only $c$ and $\sigma_{\rm core}$;
it does not supply their line shapes. In the standardized panels, a
piecewise-linear CDF through the native 0.1~MeV histogram edges is
integrated into bins of width $\Delta u=0.02$. This is empirical
rebinning, with uniform density within each original bin, rather than
an analytic signal parameterization. The dashed black curve is a
separate candidate: the equal-mass mean of these standardized shapes.

\fig{.53}{../figures/MC_points_vs_figure12.pdf}{Direct MC-point overlays
for the outlying 60~MeV sample and a representative 160~MeV sample.
Top: full-normalized native empirical curves and counts summed into
2~MeV bins. Bottom: exactly the standardized empirical curves of
Figure 12 and direct sums into 0.2~MeV bins, both conditioned
on $|u|<2$. Horizontal bars show bin widths; vertical bars show
$\sqrt{\sum w^2}$ with the display normalization held fixed. The dashed
common shape is shown separately. Points are bin averages; the curve
must be integrated over the same bins for an exact comparison.}

The native curves agree with the MC points by construction: the
largest difference between a direct bin sum and the corresponding
CDF integral is below $10^{-9}$ candidates. This verifies the plotted
representation, not an independent goodness of fit; no fit probability
is assigned to that identity. The common-shape candidate need not
agree with each sample, and the held-out tests above assess that
separate question. The plot and all underlying counts are saved.

\clearpage
\section{The 60 MeV outlier and the meaning of a window in $u$}
At 60~MeV, $32.4\%$ of the selected histogram is inside $|u|<2$,
compared with $60.7$--$71.5\%$ at the other nine masses. Around the
current centered $\pm2.25\sigma_{\rm nominal}$ window, $5.4\%$ is
on the left, $55.1\%$ on the right and $39.5\%$ inside. Thus the
full-normalized core is small because much of this selected-candidate
distribution lies elsewhere. The core-conditioned panel removes that
fraction information and still shows some shape variation.

The 60~MeV files contain 243,668 selected rows. Their event cutflows
record 1,542,086 readout events and the same 243,668 events at the final
selection stage, a conditional retention of $15.8\%$. Other masses
give $34.5$--$61.6\%$ (Table~\ref{tab:MCunits}). This supports lower
retention from the recorded readout stage, but neither that stage nor
the unit-normalized Figure 12 establishes generated-signal acceptance.
A certified generated-event denominator and signal-daughter association
are still needed; the selected candidate distribution is not certified
as a pure signal response. ``Within the study range'' does not mean
high detector acceptance.

The existing core-centered window and a proposed fixed-$u$ window are
different because $u$ uses the fitted core width:
\begin{equation}
 |m_{ee}-c|<2.25\sigma_{\rm nominal}
 \quad\Longleftrightarrow\quad
 |u|<2.25\frac{\sigma_{\rm nominal}}{\sigma_{\rm core}}
 =2.920\text{--}3.174.
\end{equation}
Consequently $|u|<2.25$ would narrow the physical window by
$23.0$--$29.1\%$. At 60~MeV, $c=58.792$~MeV,
$\sigma_{\rm core}=1.550$~MeV and $\sigma_{\rm nominal}=2.075$~MeV:
the half-width changes from 4.669 to 3.488~MeV and the retained
selected-MC fraction falls from $39.5\%$ to $34.5\%$.

\begin{table}[H]\centering\scriptsize
\begin{tabular}{rrrrrrr}
\toprule
$m_0$ & $\sigma_{core}$ & $\sigma_{nom}$ & Old half-width ($u$) & Old MC (\%) & $\pm2.25u$ MC (\%) & Selected/readout (\%) \\
\midrule
60 & 1.550 & 2.075 & 3.012 & 39.5 & 34.5 & 15.8 \\
80 & 1.637 & 2.288 & 3.145 & 70.4 & 63.6 & 34.5 \\
100 & 1.853 & 2.570 & 3.120 & 77.6 & 71.1 & 47.5 \\
120 & 2.112 & 2.920 & 3.111 & 80.5 & 74.0 & 55.3 \\
140 & 2.392 & 3.339 & 3.140 & 81.5 & 74.6 & 59.5 \\
160 & 2.713 & 3.827 & 3.174 & 81.3 & 73.9 & 61.3 \\
180 & 3.120 & 4.383 & 3.161 & 80.7 & 72.9 & 61.6 \\
200 & 3.577 & 5.008 & 3.150 & 79.2 & 71.0 & 60.0 \\
220 & 4.233 & 5.702 & 3.031 & 77.1 & 69.4 & 57.0 \\
240 & 4.981 & 6.465 & 2.920 & 74.5 & 67.3 & 52.8 \\
\bottomrule
\end{tabular}

\caption{Widths in MeV and continuous-window probabilities from the
stored 0--400~MeV MC histograms. ``Old'' means the existing
core-centered $\pm2.25\sigma_{\rm nominal}$ window. The last column
is conditional readout-stage retention, not detector acceptance.
Actual observed-bin masks use bin centers and can differ slightly
from these continuous fractions.}\label{tab:MCunits}
\end{table}

A fixed-$u$ prescription is a coherent MC-based candidate, but the
coefficient must be declared as a new choice. It is not a change of
units that preserves the old blind region. Nor does a non-Gaussian
empirical shape inherit a Gaussian containment fraction from its fitted
core width. This clarification does not change the primary mask or
recompute observed limits and local probabilities. Testing a new mask
would require the matching GP training exclusion and signal-recovery
checks, as for the preceding width diagnostics.

\end{document}
